Considera la funció
\[
f(x)=x^3+x^2+x-1 .
\]

\begin{apartats}

\apartat[1,25]{1}
Enuncia el teorema de Bolzano i demostra que l'equació $f(x)=0$ té almenys una solució
a l'interval $[0,1]$.

\begin{solucio}
\emph{Teorema de Bolzano.} Si $f$ és contínua a $[a,b]$ i $f(a)$ i $f(b)$ tenen signes
oposats, aleshores existeix almenys un $c\in(a,b)$ tal que $f(c)=0$.\\
$f$ és polinòmica, per tant contínua a $\mathbb{R}$ i en particular a $[0,1]$.
Com que $f(0)=-1<0$ i $f(1)=2>0$, hi ha almenys una arrel a $(0,1)$.
\end{solucio}

\apartat[1,25]{0,75}
Aplicant el mètode de la bisecció, troba un interval de longitud $0{,}25$ que contingui
una solució de l'equació. Justifica cada pas.

\begin{solucio}
$f\!\left(\tfrac12\right)=\tfrac18+\tfrac14+\tfrac12-1=-\tfrac18<0$. Com que
$f(1)>0$, l'arrel és a $\left(\tfrac12,1\right)$, de longitud $\tfrac12$.\\
$f\!\left(\tfrac34\right)=\tfrac{27}{64}+\tfrac{9}{16}+\tfrac34-1=\tfrac{47}{64}>0$.
Com que $f\!\left(\tfrac12\right)<0$, l'arrel és a
$\left(\tfrac12,\tfrac34\right)$, de longitud $\tfrac14=0{,}25$.
\end{solucio}

\begin{nomesllarg}
\apartat{0,75}
Demostra que les gràfiques de les funcions $y=e^{x}$ i $y=3-x$ es tallen en algun punt
d'abscissa $x\in(0,1)$.

\begin{solucio}
Es tallen on $e^x=3-x$, és a dir on $k(x)=e^x+x-3$ s'anul·la. $k$ és contínua a
$\mathbb{R}$ per ser suma de funcions contínues, i
$k(0)=1-3=-2<0$, $k(1)=e-2\approx0{,}72>0$.
Per Bolzano existeix $c\in(0,1)$ amb $k(c)=0$, que és l'abscissa del punt de tall.
\end{solucio}
\end{nomesllarg}

\end{apartats}
